\subsubsection{中医方剂补遗：五子衍宗丸、左归丸与蛇床子}

前文列出的常用方剂（金匮肾气丸、右归丸、赞育丹等）之外，男科临床与民间使用频率很高的还有几个，补记于此：

\begin{itemize}
\item \textbf{五子衍宗丸}（枸杞子、菟丝子、覆盆子、五味子、车前子）：被称为"古今种子第一方"，是中医治疗男性不育的代表方，对应西医少精子症、弱精子症的辅助调理。现代研究以小样本临床观察与动物实验为主，证据等级有限；定位应是生活方式干预与西医病因治疗之外的辅助，而非替代；
\item \textbf{左归丸}：张景岳"阴中求阳"思路的代表，纯补真阴，适用于肾阴虚为主者（潮热、盗汗、五心烦热、舌红少苔）。与右归丸（温肾阳）恰成一对——"左阴右阳"的区分正是中医辨证选药的核心：用反了方向，等于南辕北辙，这也是"补肾药越补越燥"的常见原因之一；
\item \textbf{蛇床子}：外用温肾壮阳的代表药，古代外治方中用于阳痿与阴部湿痒。现代药理提示蛇床子素有一定的平滑肌松弛与抗炎作用；需注意其内服有毒性报道，应严格遵医嘱，不建议自行泡水或泡酒服用；
\item \textbf{知柏地黄丸、归芍地黄丸}：六味地黄丸的衍生方，分别适合阴虚火旺（潮热明显）与肝肾不足兼血虚者；早泄辨证属阴虚火旺型时，知柏地黄丸比金匮肾气丸类温阳药更常用。
\end{itemize}

\textbf{共同提醒}：中成药同样讲究辨证与疗程，"补肾药当保健品长期吃"是门诊常见的错误用法——一是体质不合则助火或碍胃，二是可能延误器质性病因（血管、激素、心理）的排查。购买时认准药品批准文号；对"三天见效"的产品保持警惕，其风险见下一节。
